\documentclass[es,boletin]{uah}

\tema{6}
\titulo{Transmisión digital paso banda}{Bandpass digital transmission}
%
\begin{document}

\titulacion{Grados TIC}
\asignatura{Teoría de la Comunicación}{Communication Theory}
\departamento{Teoría de la Señal y Comunicaciones}
\curso{2024/2025} % Do not show year

\maketitle

%%%%%%%%%%%%%%%%%%%%%%%%%%%%%%%%%%%%%%%%%%%%%%%%%%%%%%%%
%%%%%%%%%%%%%%%%%%%%%%%%%%%%%%%%%%%%%%%%%%%%%%%%%%%%%%%%
\section{Problemas básicos}
%%%%%%%%%%%%%%%%%%%%%%%%%%%%%%%%%%%%%%%%%%%%%%%%%%%%%%%%
%%%%%%%%%%%%%%%%%%%%%%%%%%%%%%%%%%%%%%%%%%%%%%%%%%%%%%%%
Este primer bloque de problemas son problemas extraídos de la bibliografía de la asignatura, y consisten en algunos cálculos básicos que es necesario dominar.




\Problema{

	\cite{Sklar} Calcule el valor esperado del número de errores de bit cometidos en un día por el receptor BPSK coherente especificado más abajo, cuando opera continuamente. 
	
	La velocidad de datos es de $5000 bits/s$. Las señales digitales de entrada son $s_1(t) = A\cdot cos(\omega_p t)$ y $s_2(t)=-A \cdot cos(\omega_p t)$, donde $A = 1 mV$ y la densidad espectral unilateral de potencia de ruido es $N_0 = 10^{-11} W/Hz$. 

}
{
$2338$ bits
}		
{

	\cite{Sklar} Calculate the expected value of the number of error bits during a day for the coherent BPSK receiver described below, under continuous operation. The data rate is $5000 bits/s$. The input digital signals are $s_1(t) = A\cdot cos(\omega_p t)$ and $s_2(t)=-A \cdot cos(\omega_p t)$, where $A = 1 mV$, and the unilateral noise power spectral density is $N_0 = 10^{-11} W/Hz$. 

}
{

$2338$ bits
}		





\Problema{

	\cite{Sklar} Un sistema BPSK coherente que opera continuamente comete errores a razón de $100$ errores al día como promedio. La velocidad de los datos es de $1000 bits/s$. La densidad espectral unilateral de potencia de ruido es $N_0 = 10^{-10} W/Hz$.

\begin{enumerate}
	\item Si el sistema es ergódico, ¿cual es la probabilidad media de error?
	\item Si el valor de la potencia media recibida se ajusta a $10^{-6} W$, ¿será suficiente esta potencia para mantener la probabilidad de error calculada en el apartado a)?
\end{enumerate} 

}
{

\begin{enumerate}
	\item $1.16 \cdot 10^{-6}$
	\item No
\end{enumerate}
}
{


	\cite{Sklar} A coherent BPSK system operating continuously produces errors at an average rate of $100$ errors per day. The data rate is $1000 bits/s$. The unilateral noise power spectral density is $N_0 = 10^{-10} W/Hz$. 


\begin{enumerate}
	\item If the system is ergodic, which is the average error probability?
	\item If the average received power is adjusted to $10^{-6}$ W, would this value be enough to keep the error probability calculated in a)?
\end{enumerate} 

}
{

\begin{enumerate}
	\item $1.16 \cdot 10^{-6}$
	\item No
\end{enumerate}
}


\Problema{

\cite{Haykin} La componente de señal de un sistema PSK coherente está definida por la expresión

\begin{displaymath}
	s(t) =A_c k sen(\omega_p t) \pm A_c \sqrt{1-k^2} cos(\omega_p t)
\end{displaymath}


donde $0 \leq t < T_b$, y el signo más corresponde al símbolo $1$ y el signo menos corresponde al símbolo $0$. El primer término del segundo miembro de la igualdad representa una componente de portadora, incluida para facilitar la sincronización entre el receptor y el transmisor. Se pide:

\begin{enumerate}
	\item Dibujar la constelación de las señales descritas. ¿Qué observaciones pueden hacerse acerca de este diagrama?
	\item Mostrar que en presencia de ruido gaussiano aditivo de media nula y densidad espectral de potencia $N_0/2$, la probabilidad media de error es:
	\begin{displaymath}
		P_e = Q \left ( \sqrt{\frac{2E_b}{N_0} (1-k^2)} \right )
	\end{displaymath}
	con $E_b = \frac{1}{2} A_c^2 T_b$
	\item Suponga que el $10\%$ de la potencia transmitida está ubicada en la componente de portadora. Determinar el valor de $E_b/N_0$ necesario para obtener una probabilidad de error de $10^{-4}$.
	\item Compare este valor de $E_b/N_0$ con el requerido en un sistema PSK convencional con la misma probabilidad de error.
\end{enumerate}

}
{
\begin{enumerate}
	\item Constelación de una PSK
	\item Demostración
	\item $\frac{E_b}{N_0} = 8.02$
	\item $\frac{E_b}{N_0} = 7.22$
\end{enumerate}

}
{


	\cite{Haykin} The signal component of a coherent PSK system is defined by the expression

\begin{displaymath}
	s(t) =A_c k sen(\omega_p t) \pm A_c \sqrt{1-k^2} cos(\omega_p t)
\end{displaymath}

where $0 \leq t < T_b$, and the plus sign corresponds to the $1$ symbol, and the minus sign corresponds to the $0$ one. The first term on the right hand side of the equation represents a carrier component, included to improve the synchronization between transmitter and receiver. Solve this:

\begin{enumerate}
	\item Plot the constellation of the signals described; what can be said about this diagram?
	\item Show that, in presence of zero-mean additive white Gaussian noise with power spectral density $N_0/2$, the average error probability is
	\begin{displaymath}
		P_e = Q \left ( \sqrt{\frac{2E_b}{N_0} (1-k^2)} \right )
	\end{displaymath}
	with $E_b = \frac{1}{2} A_c^2 T_b$
	\item Assume that $10\%$ of the transmitted power is located in the carrier component. Determine the value of $E_b/N_0$ required to obtain an error probability of $10^{-4}$.
	\item Compare this $E_b/N_0$ value with the one required in a conventional PSK system with the same error probability.
\end{enumerate}

}
{
\begin{enumerate}
	\item PSK constellation
	\item Demonstration
	\item $\frac{E_b}{N_0} = 8.02$
	\item $\frac{E_b}{N_0} = 7.22$
\end{enumerate}

}


\Problema{

	\cite{Haykin} Se desea comparar dos sistemas paso-banda de transmisión de datos. Uno de ellos utiliza 16-PSK y el otro 16-QAM. Ambos sistemas deben obtener una probabilidad media de error de símbolo igual a $10^{-3}$. Compare los requisitos de relación señal-ruido de estos dos sistemas.
}
{

$\Delta \left ( \frac{E_s}{N_0} \right ) = 3.68 dB$
}
{

	\cite{Haykin} We want to compare two data transmission bandpass systems. One of them employs 16-PSK, the other, 16-QAM. Both systems have to provide an average symbol error probability of $10^{-3}$. Compare the signal-to-noise requirements of said systems.
}
{

$\Delta \left ( \frac{E_s}{N_0} \right ) = 3.68 dB$
}

\Problema{

\cite{Sklar} Si el criterio de rendimiento de un sistema es la probabilidad de error de bit, ¿Cuál de los dos esquemas de modulación siguientes será el elegido para operar en un canal AWGN? Muestre los cálculos.

\begin{enumerate}
	\item FSK binaria ortogonal no coherente con $E_b/N_0 = 13 dB$.
	\item PSK binaria coherente con $E_b/N_0 = 8 dB$.
\end{enumerate}

}
{

FSK binaria ortogonal no coherente
}
{

	\cite{Sklar} If the performance criterion of a system is the bit error probability, which one of the following modulation schemes would be chosen to operate in an AWGN channel? Show the calculations.

\begin{enumerate}
	\item Non-coherent binary orthogonal FSK with $E_b/N_0 = 13 dB$.
	\item Coherent binary PSK with $E_b/N_0 = 8 dB$.
\end{enumerate}

}
{

Non-coherent binary orthogonal FSK
}



\Problema{

\cite{Haykin} 	Un sistema transmite datos binarios a la velocidad de $2,5 \cdot 10^6$ bps. En transcurso de la transmisión se añade a la señal un ruido blanco gaussiano de media nula y densidad espectral de potencia $N_0/2=10^(-20)$ W/Hz. En ausencia de ruido, la amplitud de la señal sinusoidal recibida para el dígito 1 ó 0 es $1 \mu V$. Determine la probabilidad media de error de símbolo para las configuraciones del sistema siguientes:

\begin{enumerate}
	\item FSK binaria coherente.
	\item MSK coherente.
	\item FSK binaria no coherente
\end{enumerate}

}
{

\begin{enumerate}
	\item $9,67 \cdot 10^{-4}$
	\item $5,41 \cdot 10^{-6}$
	\item $3,36 \cdot 10^{-3}$
\end{enumerate}
}
{

\cite{Haykin} A system transmits binary data at a rate of $2.5 \cdot 10^6$ bps. During transmission, zero-mean additive white Gaussian noise with power spectral density $N_0/2 = 10^{-20}$ W/Hz is added to the signal. In the absence of noise, the amplitude of the received sinusoidal signal for digit 1 or 0 is $1\,\mu V$. Determine the average symbol error probability for the following system configurations:

\begin{enumerate}
	\item Coherent binary FSK.
	\item Coherent MSK.
	\item Non-coherent binary FSK.
\end{enumerate}

}
{

\begin{enumerate}
	\item $9,67 \cdot 10^{-4}$
	\item $5,41 \cdot 10^{-6}$
	\item $3,36 \cdot 10^{-3}$
\end{enumerate}
}

\Problema{
		(Examen mayo de 2016) Un sistema de comunicaciones digitales emplea un esquema de modulación FSK para transmitir dos señales ortogonales y equiprobables $s_1(t)$ y $s_2(t)$ a la velocidad de $2,5\cdot 10^6$ bps. El sistema está caracterizado por $E_b/N_0 = 12$ dB.

$s_i (t) = \begin{cases} \sqrt{\frac{2E_b}{T_b}} \cos(\omega_i t), & 0 \leq t < T_b \\ 0, & \text{c.c.} \end{cases}, \quad i = 1,2.$

\begin{enumerate}
	\item Calcule la probabilidad media de error para los esquemas de modulación FSK con receptor coherente y receptor no coherente.
	\item ¿Cuál de los esquemas anteriores es más adecuado para un canal modelado con ruido AWGN? Justifique su respuesta. Para el resto del ejercicio tenga en cuenta solo el esquema seleccionado en este apartado.
	\item Calcule el ancho de banda mínimo del esquema seleccionado en el apartado b).
	\item Para dicho esquema, determine una base ortonormal que permita representar $s_1(t)$ y $s_2(t)$. Represente la constelación en el espacio de señal indicando las coordenadas de los vectores señal.
	\item Si compara la modulación anterior con una BPSK coherente con la misma $P_e$, ¿Cuál de los dos exigirá mayor relación $E_b/N_0$? Calcule el incremento de $E_b/N_0$ expresado en dB.
\end{enumerate}	
}
{
\begin{enumerate}
	\item $P_e^{coh} = 3,43 \cdot 10^{-5}$, $P_e^{no-coh} = 1.81 \cdot 10^{-4}$
	\item FSK con receptor coherente
	\item $B_{min} = 6,25$ MHz
	\item Base: $s_1(t)$ y $s_2(t)$ normalizadas. 
	\item La FSK. Incremento de $E_b/N_0$ de $3$ dB.
\end{enumerate}
}	
{
(May 2016 exam) A digital communications system uses an FSK modulation scheme to transmit two orthogonal and equiprobable signals, $s_1(t)$ and $s_2(t)$, at a rate of $2.5\cdot 10^6$ bps. The system is characterized by $E_b/N_0 = 12$ dB.

$s_i (t) = \begin{cases} \sqrt{\frac{2E_b}{T_b}} \cos(\omega_i t), & 0 \leq t < T_b \\ 0, & \text{otherwise} \end{cases}, \quad i = 1,2.$

\begin{enumerate}
	\item Calculate the average error probability for FSK modulation schemes with coherent and non-coherent receivers.
	\item Which of the previous schemes is more suitable for a channel modeled with AWGN noise? Justify your answer. For the rest of the exercise, consider only the scheme selected in this section.
	\item Calculate the minimum bandwidth of the scheme selected in part b).
	\item For that scheme, determine an orthonormal basis that allows representing $s_1(t)$ and $s_2(t)$. Plot the constellation in signal space indicating the coordinates of the signal vectors.
	\item If you compare the previous modulation with a coherent BPSK with the same $P_e$, which one will require a higher $E_b/N_0$ ratio? Compute the increase in $E_b/N_0$ expressed in dB.
\end{enumerate}	
}
{
\begin{enumerate}
	\item $P_e^{coh} = 3.43 \cdot 10^{-5}$, $P_e^{no-coh} = 1.81 \cdot 10^{-4}$
	\item FSK with coherent receiver
	\item $B_{min} = 6.25$ MHz
	\item Basis: normalized $s_1(t)$ and $s_2(t)$. 
	\item FSK. $E_b/N_0$ increase of $3$ dB.
\end{enumerate}
}


\begin{thebibliography}{100}
	\bibitem[Haykin2001]{Haykin}  Simon Haykin. Communication Systems, 4th Ed. John Wiley and Sons, 2001.
	\bibitem[Sklar2001]{Sklar} Bernard Sklar. Digital Communications, 2nd Ed. Prentice Hall, 2001.
	\end{thebibliography}

\end{document}


